\documentclass[11pt]{article}
\usepackage{preamble}
\setlength{\parskip}{4pt}

\begin{document}

\daybanner{AP Statistics \textperiodcentered\ Unit 2 \textperiodcentered\ Topic 2.12 \textperiodcentered\ B: 2026-11-12 \textperiodcentered\ E: 2026-12-07 \textperiodcentered\ TEACHER KEY}{One Sample Far From the Others}

\sectionbanner{CED TETHER}

{\small
\begin{itemize}[leftmargin=1.5em,itemsep=2pt,topsep=2pt]
  \item \textbf{VAR-1.G} Identify questions suggested by variation in statistics for samples collected from the same population. \textbf{[Skill 1.A]}
  \item \textbf{VAR-1.G.1} Variation in statistics for samples taken from the same population may be random or not.
\end{itemize}
}
\textbf{Essential question:} When is an unusual sample statistic evidence of an error, and when is it ordinary sampling variability?\par
\textbf{DOK-3 skill:} 4.B \textperiodcentered\ \textbf{FRQ pattern:} {\small\texttt{sampling-\allowbreak{}variability-\allowbreak{}is-\allowbreak{}not-\allowbreak{}error}}\par
\textbf{DOK rationale:} Part (c) weighs an unusual sample mean against a graph and distinguishes a fixed population parameter from random variation to decide whether the evidence establishes a collection error.\par

\frameworkphaseheader{Do Now (first take) $\rightarrow$ Explore (video + follow-along) $\rightarrow$ Exit (finish + turn in)}{1 $\rightarrow$ 3}{45}{%
\begin{itemize}[leftmargin=1.5em,itemsep=2pt,topsep=2pt]
  \item Board slide up at the bell. Ask students to classify 520 before defining sampling variability.
  \item Begin the combined 5.1--5.2 video follow-along at minute 5.
  \item Return to the board slide at minute 33. Circulate and separate unusualness from evidence of a collection error.
  \item Collect at the door. Score part (c) only, E/P/I.
\end{itemize}
}{%
\begin{itemize}[leftmargin=1.5em,itemsep=2pt,topsep=2pt]
  \item Choose error or sampling variability and cite one feature of the dotplot.
  \item Complete the follow-along about statistics from repeated samples and predictable sample-to-sample variation.
  \item Finish (a)--(c), revise the first take in the exit line, and turn in the sheet.
\end{itemize}
}{%
\begin{itemize}[leftmargin=1.5em,itemsep=2pt,topsep=2pt]
  \item Which number belongs to the whole population, and which belongs to Elena's ten packages?
  \item Do the other 19 sample means match one another exactly?
  \item What evidence about the sampling procedure would you need before calling 520 an error?
  \item Would more repetitions or more packages per sample shrink the distribution itself?
\end{itemize}
}{Solo teacher. Circulate ~5 pairs. Require students to point to both ordinary spread and the isolated value; do not equate unusual with impossible or incorrect.}

\begin{callout}{\IconWarn\ WATCH FOR}{calloutred}
\begin{itemize}[leftmargin=1.5em,itemsep=2pt,topsep=2pt]
  \item Treating the population mean as the required result of every random sample.
  \item Calling the 12-gram gap proof of a recording mistake.
  \item Saying more simulated repetitions shrink variability when each sample remains n equals 10.
\end{itemize}
\end{callout}

\sectionbanner{SCORING PART (c) --- E / P / I}

\textbf{Expected elements}
\begin{itemize}[leftmargin=1.5em,itemsep=2pt,topsep=2pt]
  \item \textbf{judgment-1} (required) --- Does not infer a collection error solely from an unusual result.
  \item \textbf{judgment-2} (required) --- Uses the 492--520 spread or 12-gram gap as evidence of unusualness.
  \item \textbf{judgment-3} (required) --- Explains fixed population mean versus sample means that vary by chance.
\end{itemize}

{\small\renewcommand{\arraystretch}{1.25}
\noindent\begin{tabular}{|p{0.5in}|p{5.9in}|}
\hline
\textbf{E} & Meets all three checks with correct contextual reasoning; equivalent wording is acceptable. \\ \hline
\textbf{P} & Shows at least one substantive correct check, but has an omission or error that prevents E. \\ \hline
\textbf{I} & Shows none of the three checks with substantive correct reasoning; a choice alone is insufficient. \\ \hline
\end{tabular}}

\textbf{Common mistakes}
\begin{itemize}[leftmargin=1.5em,itemsep=2pt,topsep=2pt]
  \item Treating the population mean of 500 as a requirement for every random sample mean
  \item Calling 520 a data-entry error without evidence about how the sample was generated
  \item Confusing the spread within one sample with the spread of means across samples
  \item Increasing only the number of simulated samples rather than the size of each sample
\end{itemize}

\clearpage
\focusbanner{TODAY'S PROBLEM --- annotated}

Suppose a class uses a simulator containing a very large population of package weights with population mean $\mu=500$ grams. Each of 20 students independently takes a random sample of $n=10$ packages and records the sample mean $\bar{x}$. The dotplot shows the 20 sample means. Elena's sample mean is 520 grams. She says, ``My sample must be wrong because its mean is far from everyone else's.'' Another student says, ``A random sample can be unusual without being collected incorrectly.''\par
\begin{center}
\begin{tikzpicture}
\begin{axis}[
  width=0.61\linewidth,
  height=1.15in,
  ymin=0, ymax=4, ytick=\empty, axis y line=none, axis x line=bottom,
  xlabel={Sample mean package weight (grams), n=10}, tick label style={font=\small}, label style={font=\small}
]
\addplot[only marks, mark=*, mark size=1.7pt, color=dayblue] coordinates {
  (492,1)
  (494,1)
  (495,1)
  (496,1)
  (497,1)
  (498,1)
  (499,1)
  (499,2)
  (500,1)
  (500,2)
  (500,3)
  (501,1)
  (501,2)
  (502,1)
  (503,1)
  (504,1)
  (505,1)
  (506,1)
  (508,1)
  (520,1)
};
\end{axis}
\end{tikzpicture}
\end{center}
\begin{firsttakebox}{FIRST TAKE (5 min)}
Does the dot at 520 prove Elena made a mistake? Give one reason from the picture.\par
\answer{Not graded. Expect students to distrust the isolated value; the finish should preserve unusualness while removing the unsupported error claim.}
\end{firsttakebox}

\textit{Rules callout ``PARAMETER FIXED, STATISTIC VARIABLE (keep on the page)'' is printed on the student sheet.}\par\medskip

\dokbadge{1}~\textbf{(a)} Identify the population parameter and Elena's sample statistic, including symbols, values, and units.\par
\answer{The population parameter is the population mean $\mu=500$ grams. Elena's sample statistic is her sample mean $\bar{x}=520$ grams from $n=10$ packages.}

\dokbadge{2}~\textbf{(b)} Describe the center and overall range of the sample means, and the gap next to 520 grams. Two sentences are enough.\par
\answer{The sample means are centered near 500 grams and range from 492 to 520 grams, a range of 28 grams. There is a 12-gram gap between 508 and 520.}

\dokbadge{3}~\textbf{(c)} In 3--4 sentences, decide whether the dotplot proves Elena's sample was collected incorrectly. Use the spread or gap as evidence. Explain how the population mean can stay fixed while correctly collected sample means differ.\par
\answer{The dotplot does not prove a collection error. The sample means range from 492 to 520 grams, with a 12-gram gap before 520, so Elena has an unusual result. The population mean $\mu=500$ grams stays fixed, but each sample contains different packages and can have a different mean by chance. We would need evidence about collection to call her sample wrong.}

\begin{summaryexitbox}{EXIT}
One thing the video changed about my first take: \hrulefill
\end{summaryexitbox}

\end{document}
